\begin{tikzpicture}[x=0.95cm,y=0.78cm,line cap=round,line join=round]
\coordinate (A) at (-1.6,0);
\coordinate (C) at (2.4,0);
\coordinate (B) at (0.4,-1.65);
\coordinate (S) at (-1.6,3.1);
\coordinate (D) at (-0.8,0);
\coordinate (E) at (0.4,1.55);

% 依原卷中点、垂直平分线及棱锥的关系重绘；不表达额外长度比例。
\draw[thick] (A)--(S)--(C)--(B)--cycle;
\draw[thick] (S)--(B);
\draw[dashed] (A)--(C) (D)--(B);
\draw[dashed,thick] (D)--(E);
\draw[thick] (S)--(E);
\draw[thick] (E)--(C);

% SC 的等长刻痕。D、E 的投影位置保留原卷的垂直平分线关系。
\draw (-0.71,2.185) -- (-0.49,2.465);
\draw (1.29,0.635) -- (1.51,0.915);
% DE 与 SC 的直角标记，按两边方向构造小方角。
\draw (0.266,1.376) -- (0.44,1.242) -- (0.574,1.416);

\node[above left] at (S) {$S$};
\node[left] at (A) {$A$};
\node[below] at (B) {$B$};
\node[right] at (C) {$C$};
\node[below] at (D) {$D$};
\node[above right] at (E) {$E$};
\node[font=\small,align=center] at (0.4,-2.35)
  {依原卷关系重绘的示意图；非比例图。};
\end{tikzpicture}
